\begin{proof}[Proof of \cref{obj:lem:moment-matched-reciprocal-priors}]
\begin{enumerate}
\item Fix \(K\ge2\), and define
\[
H=K^2,\qquad
\mathcal I_\star=[1,H],\qquad
f_\star:\mathcal I_\star\to\mathbb R,\quad f_\star(z)=z^{-1}.
\]
Since \(H>1\) and every point of \(\mathcal I_\star\) is positive, \(f_\star\) is continuous. For a continuous function \(\varphi:\mathcal I_\star\to\mathbb R\), write
\[
\|\varphi\|_\infty=\max_{z\in\mathcal I_\star}|\varphi(z)|,
\qquad
E_K=\inf_{\substack{\mathfrak p\in\mathbb R[z]\\ \deg\mathfrak p\le K}}
\|f_\star-\mathfrak p\|_\infty.
\]
In particular, \(\|f_\star\|_\infty=1\).

To obtain a best approximant, consider the finite-dimensional space
\[
\mathcal V_K
=\{\mathfrak p|_{\mathcal I_\star}:
  \mathfrak p\in\mathbb R[z],\ \deg\mathfrak p\le K\},
\qquad
\mathcal B_\star
=\{\varphi\in\mathcal V_K:\|\varphi\|_\infty\le3\}.
\]
The closed ball \(\mathcal B_\star\) is compact. The continuous error function therefore attains its minimum there at the restriction of a polynomial \(p_\star\) of degree at most \(K\). Comparison with the zero polynomial gives
\[
\|f_\star-p_\star\|_\infty\le1.
\]
For every \(\varphi\in\mathcal V_K\setminus\mathcal B_\star\), the reverse triangle inequality gives
\[
\|f_\star-\varphi\|_\infty
\ge\|\varphi\|_\infty-\|f_\star\|_\infty
>3-1=2
\ge\|f_\star-p_\star\|_\infty.
\]
Thus this minimum is global:
\[
\deg p_\star\le K,\qquad
\|f_\star-p_\star\|_\infty=E_K\ge0.
\]
% lean: Causalean.Mathlib.Analysis.FinitePolynomialAlternationDuality.exists_bestPolynomial

\item Define the residual
\[
e_\star(z)=f_\star(z)-p_\star(z)
\qquad(z\in\mathcal I_\star).
\]
We establish the existence of points and an orientation satisfying
\[
1\le z_0<z_1<\cdots<z_{K+1}\le H,\qquad
\eta_\star\in\{-1,1\},\qquad
e_\star(z_\iota)=\eta_\star(-1)^\iota E_K
\quad(0\le\iota\le K+1).
\]
If \(E_K=0\), the residual vanishes throughout the interval, and the choices
\[
z_\iota=1+\frac{\iota(H-1)}{K+1}
\quad(0\le\iota\le K+1),
\qquad
\eta_\star=1
\]
have these properties.

Suppose now that \(E_K>0\). Define the compact extremal sets and the residual sign by
\[
\begin{aligned}
\mathcal Z_\star
&=\{z\in\mathcal I_\star:|e_\star(z)|=E_K\},\\
\mathcal Z_\star^+
&=\{z\in\mathcal Z_\star:e_\star(z)=E_K\},\\
\mathcal Z_\star^-
&=\{z\in\mathcal Z_\star:e_\star(z)=-E_K\},\\
\chi_\star(z)
&=
\begin{cases}
1,&e_\star(z)>0,\\
0,&e_\star(z)=0,\\
-1,&e_\star(z)<0
\end{cases}
\qquad(z\in\mathcal I_\star).
\end{aligned}
\]
The set \(\mathcal Z_\star\) is nonempty because the maximum residual magnitude is attained.

Assume, for a contradiction, that the required \(K+2\) alternating extrema do not exist. We first construct a polynomial \(p_{\rm sign}\) such that
\[
\deg p_{\rm sign}\le K,\qquad
e_\star(z)p_{\rm sign}(z)>0
\quad(z\in\mathcal Z_\star).
\]
If \(\mathcal Z_\star^+=\varnothing\), take \(p_{\rm sign}(z)=-1\); if \(\mathcal Z_\star^-=\varnothing\), take \(p_{\rm sign}(z)=1\). These choices give
\[
e_\star(z)p_{\rm sign}(z)=E_K>0
\quad(z\in\mathcal Z_\star).
\]

When both sign sets are nonempty, their disjointness and compactness give a separation \(\delta_{\rm sep}>0\) with
\[
\delta_{\rm sep}<|z-z'|
\quad(z\in\mathcal Z_\star^+,\ z'\in\mathcal Z_\star^-).
\]
Set
\[
\varepsilon_{\rm sep}=\frac{\delta_{\rm sep}}5.
\]
Choose finitely many distinct centers in \(\mathcal Z_\star\), ordered as
\[
c_0^{\rm ctr}<\cdots<c_{n_{\rm ctr}-1}^{\rm ctr},
\qquad n_{\rm ctr}\ge1,
\]
whose \(\varepsilon_{\rm sep}\)-neighborhoods cover \(\mathcal Z_\star\). Thus
\[
\forall z\in\mathcal Z_\star\ \exists\,\iota\in\{0,\ldots,n_{\rm ctr}-1\}:
\quad |z-c_\iota^{\rm ctr}|<\varepsilon_{\rm sep}.
\]
Define the consecutive sign-change indices by
\[
\mathcal D_{\rm ctr}
=\{\iota\in\{0,\ldots,n_{\rm ctr}-2\}:
  \chi_\star(c_\iota^{\rm ctr})
  \ne\chi_\star(c_{\iota+1}^{\rm ctr})\}.
\]
Selecting one center from each successive sign block produces an alternating subsequence of length \(|\mathcal D_{\rm ctr}|+1\). The assumed absence of \(K+2\) alternating extrema therefore implies
\[
|\mathcal D_{\rm ctr}|\le K.
\]
Define
\[
p_{\rm sign}(z)
=\chi_\star(c_{n_{\rm ctr}-1}^{\rm ctr})
 \prod_{\iota\in\mathcal D_{\rm ctr}}
 \left(z-\frac{c_\iota^{\rm ctr}+c_{\iota+1}^{\rm ctr}}2\right),
\]
with an empty product equal to one. Its degree is \(|\mathcal D_{\rm ctr}|\le K\). Each factor changes sign exactly between consecutive blocks, so its sign at every center equals the residual sign there:
\[
\chi_\star(c_\iota^{\rm ctr})p_{\rm sign}(c_\iota^{\rm ctr})>0.
\]
Oppositely signed centers are more than \(5\varepsilon_{\rm sep}\) apart. Each midpoint appearing as a root is therefore more than \(2\varepsilon_{\rm sep}\) from every center. Consequently,
\[
\chi_\star(c_\iota^{\rm ctr})p_{\rm sign}(z)>0
\quad\text{whenever }|z-c_\iota^{\rm ctr}|<\varepsilon_{\rm sep}.
\]
An extremal point and a center within \(\varepsilon_{\rm sep}\) have the same residual sign, by the separation of \(\mathcal Z_\star^+\) and \(\mathcal Z_\star^-\). Hence
\[
e_\star(z)p_{\rm sign}(z)
=E_K\chi_\star(z)p_{\rm sign}(z)>0
\quad(z\in\mathcal Z_\star).
\]
% lean: Causalean.Mathlib.Analysis.FinitePolynomialAlternationDuality.exists_signPolynomial_of_no_alternatingExtrema

We next turn this sign agreement into a strict improvement. Define
\[
\mathcal B_{\rm bad}
=\{z\in\mathcal I_\star:e_\star(z)p_{\rm sign}(z)\le0\},
\qquad
c_{\rm cut}
=
\begin{cases}
\displaystyle
\frac{E_K+\max_{z\in\mathcal B_{\rm bad}}|e_\star(z)|}{2},
&\mathcal B_{\rm bad}\ne\varnothing,\\[6pt]
E_K/2,&\mathcal B_{\rm bad}=\varnothing.
\end{cases}
\]
When nonempty, \(\mathcal B_{\rm bad}\) is compact and contains no extremal point. Its maximum residual magnitude is therefore strictly below \(E_K\). In either case,
\[
c_{\rm cut}<E_K,\qquad
|e_\star(z)|\ge c_{\rm cut}
\ \Longrightarrow\
e_\star(z)p_{\rm sign}(z)>0.
\]
Define
\[
\begin{aligned}
\mathcal U_{\rm high}
&=\{z\in\mathcal I_\star:|e_\star(z)|\ge c_{\rm cut}\},\\
\mu_{\rm sign}
&=\min_{z\in\mathcal U_{\rm high}}
  e_\star(z)p_{\rm sign}(z)>0,\\
M_{\rm sign}
&=\|p_{\rm sign}\|_\infty>0,\\
t_\star
&=\min\left\{
\frac{\mu_{\rm sign}}{M_{\rm sign}^2},
\frac{E_K-c_{\rm cut}}{2M_{\rm sign}}
\right\}>0.
\end{aligned}
\]
The minimum defining \(\mu_{\rm sign}\) exists because \(\mathcal U_{\rm high}\) is compact and contains \(\mathcal Z_\star\). The positivity of \(M_{\rm sign}\) follows from the strict sign agreement at an extremal point.

For \(z\in\mathcal U_{\rm high}\),
\[
\begin{aligned}
\bigl(e_\star(z)-t_\star p_{\rm sign}(z)\bigr)^2
&=e_\star(z)^2
 -2t_\star e_\star(z)p_{\rm sign}(z)
 +t_\star^2p_{\rm sign}(z)^2\\
&\le E_K^2-2t_\star\mu_{\rm sign}
       +t_\star^2M_{\rm sign}^2\\
&\le E_K^2-t_\star\mu_{\rm sign}
<E_K^2.
\end{aligned}
\]
For \(z\in\mathcal I_\star\setminus\mathcal U_{\rm high}\),
\[
\begin{aligned}
|e_\star(z)-t_\star p_{\rm sign}(z)|
&\le |e_\star(z)|+t_\star M_{\rm sign}\\
&<c_{\rm cut}+\frac{E_K-c_{\rm cut}}2
=\frac{E_K+c_{\rm cut}}2<E_K.
\end{aligned}
\]
Continuity and compactness now give
\[
\|f_\star-(p_\star+t_\star p_{\rm sign})\|_\infty<E_K,
\qquad
\deg(p_\star+t_\star p_{\rm sign})\le K,
\]
contradicting the definition of \(E_K\). Thus the required alternating extrema exist.
% lean: Causalean.Mathlib.Analysis.FinitePolynomialAlternationDuality.exists_equioscillationWitness

\item Using these nodes, define the index set, nodal weights, normalization, and signed weights by
\[
\begin{aligned}
\mathcal N_K&=\{0,\ldots,K+1\},\\
b_\iota^{\rm lag}
&=\left(\prod_{\substack{\iota'\in\mathcal N_K\\\iota'\ne\iota}}
(z_\iota-z_{\iota'})\right)^{-1},\\
Z_{\rm lag}&=\sum_{\iota\in\mathcal N_K}|b_\iota^{\rm lag}|,\\
w_\iota&=\frac{b_\iota^{\rm lag}}{Z_{\rm lag}}
\qquad(\iota\in\mathcal N_K).
\end{aligned}
\]
All nodal weights are nonzero, so \(Z_{\rm lag}>0\). There are exactly \(K+1-\iota\) negative factors in the product defining \(b_\iota^{\rm lag}\). It follows that
\[
\sum_{\iota\in\mathcal N_K}|w_\iota|=1,
\qquad
w_\iota=(-1)^{K+1-\iota}|w_\iota|.
\]

For any polynomial \(\mathfrak p\) of degree at most \(K\), Lagrange interpolation gives
\[
\mathfrak p(z)
=\sum_{\iota\in\mathcal N_K}
\mathfrak p(z_\iota)b_\iota^{\rm lag}
\prod_{\substack{\iota'\in\mathcal N_K\\\iota'\ne\iota}}
(z-z_{\iota'}).
\]
Comparing coefficients of \(z^{K+1}\), whose coefficient on the left is zero, and dividing by \(Z_{\rm lag}\), yields
\[
\sum_{\iota\in\mathcal N_K}w_\iota\mathfrak p(z_\iota)=0.
\]
In particular,
\[
\sum_{\iota\in\mathcal N_K}w_\iota z_\iota^v=0
\quad(0\le v\le K),
\qquad
\sum_{\iota\in\mathcal N_K}w_\iota p_\star(z_\iota)=0.
\]
Define the signed reciprocal evaluation by
\[
d_\star=\sum_{\iota\in\mathcal N_K}w_\iota z_\iota^{-1}.
\]
Polynomial cancellation, alternation, and normalization give
\[
\begin{aligned}
d_\star
&=\sum_{\iota\in\mathcal N_K}
  w_\iota\bigl(f_\star(z_\iota)-p_\star(z_\iota)\bigr)\\
&=\sum_{\iota\in\mathcal N_K}
  (-1)^{K+1-\iota}|w_\iota|
  \eta_\star(-1)^\iota E_K\\
&=\eta_\star(-1)^{K+1}E_K
  \sum_{\iota\in\mathcal N_K}|w_\iota|
=\eta_\star(-1)^{K+1}E_K.
\end{aligned}
\]
Since \(E_K\ge0\) and \(|\eta_\star|=1\), this proves
\[
|d_\star|=E_K.
\]
% lean: Causalean.Mathlib.Analysis.FinitePolynomialAlternationDuality.exists_finiteMomentDual

\item Define two nonnegative atomic measures by
\[
\Pi^+
=2\sum_{\iota\in\mathcal N_K}\max(w_\iota,0)\,\delta_{z_\iota},
\qquad
\Pi^-
=2\sum_{\iota\in\mathcal N_K}\max(-w_\iota,0)\,\delta_{z_\iota},
\]
where the unit point mass is specified, for every Borel set \(\mathcal T\subseteq\mathbb R\), by
\[
\delta_z(\mathcal T)
=
\begin{cases}
1,&z\in\mathcal T,\\
0,&z\notin\mathcal T.
\end{cases}
\]
The moment identity at \(v=0\) gives \(\sum_\iota w_\iota=0\). Applying
\[
|w_\iota|=2\max(w_\iota,0)-w_\iota
         =2\max(-w_\iota,0)+w_\iota
\]
and summing, we obtain
\[
1=2\sum_{\iota\in\mathcal N_K}\max(w_\iota,0)
 =2\sum_{\iota\in\mathcal N_K}\max(-w_\iota,0).
\]
Thus \(\Pi^+\) and \(\Pi^-\) are probability measures.

Orient the pair according to the reciprocal evaluation:
\[
(\Pi_0,\Pi_1)
=
\begin{cases}
(\Pi^-,\Pi^+),&d_\star\ge0,\\
(\Pi^+,\Pi^-),&d_\star<0.
\end{cases}
\]
Both oriented measures are therefore probability measures. Define their common finite supporting set by
\[
\mathcal S_\star=\{z_\iota:\iota\in\mathcal N_K\}.
\]
Every atom belongs to \(\mathcal S_\star\subseteq[1,H]\), and hence
\[
\Pi^+(\mathcal S_\star)=\Pi^-(\mathcal S_\star)=1,
\qquad
\Pi_0(\mathcal S_\star)=\Pi_1(\mathcal S_\star)=1.
\]
% lean: CausalSmith.Stat.MarRareqLogfrontier.moment_matched_reciprocal_priors

\item For any real-valued function \(\varphi_{\rm test}:\mathbb R\to\mathbb R\), the finite atomic formulas and
\(\max(w_\iota,0)-\max(-w_\iota,0)=w_\iota\) give
\[
\begin{aligned}
\int\varphi_{\rm test}\,d\Pi^+
-\int\varphi_{\rm test}\,d\Pi^-
&=2\sum_{\iota\in\mathcal N_K}
 \bigl(\max(w_\iota,0)-\max(-w_\iota,0)\bigr)
 \varphi_{\rm test}(z_\iota)\\
&=2\sum_{\iota\in\mathcal N_K}
 w_\iota\varphi_{\rm test}(z_\iota).
\end{aligned}
\]
Taking \(\varphi_{\rm test}(z)=z^v\), for an integer \(0\le v\le K\), yields
\[
\int z^v\,d\Pi^+(z)-\int z^v\,d\Pi^-(z)
=2\sum_{\iota\in\mathcal N_K}w_\iota z_\iota^v=0.
\]
Either orientation therefore satisfies
\[
\int z^v\,d\Pi_0(z)=\int z^v\,d\Pi_1(z)
\qquad(0\le v\le K).
\]
% lean: Causalean.Stat.Minimax.MomentMatchedMixture.FiniteSignedMomentMarkedPoissonMixture.NormalizedFiniteSignedMomentCertificate.orientedPriors_moments_eq

For the reciprocal function, the orientation gives
\[
\int z^{-1}\,d\Pi_1(z)-\int z^{-1}\,d\Pi_0(z)
=
\begin{cases}
2d_\star,&d_\star\ge0,\\
-2d_\star,&d_\star<0
\end{cases}
=2|d_\star|=2E_K.
\]
% lean: Causalean.Stat.Minimax.MomentMatchedMixture.FiniteSignedMomentMarkedPoissonMixture.NormalizedFiniteSignedMomentCertificate.orientedPrior_target_separation

\item It remains to bound \(2E_K\). By \cref{obj:lem:reciprocal-best-approximation},
\[
E_K=\frac{K^2-1}{2K^2}
\left(\frac{K-1}{K+1}\right)^K.
\]
If \(K=2\), this gives
\[
E_2=\frac1{24},
\qquad
2E_2=\frac1{12}.
\]

Suppose \(K\ge3\). First,
\[
\frac{K^2-1}{2K^2}-\frac49
=\frac{K^2-9}{18K^2}\ge0.
\]
To bound the power term, define
\[
x_{\rm rec}=\frac1K\in(0,1/3].
\]
The logarithmic remainder estimate at order zero gives
\[
\frac12\log\frac{1+x_{\rm rec}}{1-x_{\rm rec}}
\le\frac{x_{\rm rec}}{1-x_{\rm rec}^2}.
\]
Since
\[
\frac{1+x_{\rm rec}}{1-x_{\rm rec}}
=\left(\frac{K-1}{K+1}\right)^{-1},
\qquad
8K^2\le9(K^2-1),
\]
we obtain
\[
-K\log\frac{K-1}{K+1}
\le K\frac{2x_{\rm rec}}{1-x_{\rm rec}^2}
=\frac{2K^2}{K^2-1}
\le\frac94.
\]
Exponentiating, using the positivity of \((K-1)/(K+1)\), gives
\[
e^{-9/4}\le\left(\frac{K-1}{K+1}\right)^K.
\]

For the numerical calibration, the standard numerical exponential bound and the rational exponential inequality give
\[
e<2.7182818286<\frac{11}{4},
\qquad
e^{1/4}\le\frac97.
\]
Here the latter estimate is the specialization at \(1/4\) of
\[
e^y\le\frac{2+y}{2-y}\qquad(0\le y<2).
\]
Consequently,
\[
e^{9/4}=e^2e^{1/4}
\le\left(\frac{11}{4}\right)^2\frac97
=\frac{1089}{112}
\le\frac{32}{3}.
\]
Taking reciprocals yields
\[
\frac3{32}\le e^{-9/4}
\le\left(\frac{K-1}{K+1}\right)^K.
\]
% lean: CausalSmith.Stat.MarRareqLogfrontier.reciprocal_ratio_power_lower

Combining the coefficient and power bounds,
\[
2E_K
=2\frac{K^2-1}{2K^2}
 \left(\frac{K-1}{K+1}\right)^K
\ge2\cdot\frac49\cdot\frac3{32}
=\frac1{12}.
\]
Thus, in both cases,
\[
\left|\int z^{-1}\,d\Pi_1(z)-\int z^{-1}\,d\Pi_0(z)\right|
=|2E_K|
\ge2E_K
\ge\frac1{12}.
\]
Together with the finite support and moment identities established above, this proves the result.
% lean: CausalSmith.Stat.MarRareqLogfrontier.reciprocal_best_error_lower
\end{enumerate}
\end{proof}
